\chapter{Delta-One and Dividends}\label{ch:s2:delta-one-and-dividends}

A bank that sells autocallables is left long dividends it did not want, and it sells them in dividend futures and swaps. Manley and Mueller-Glissmann
reported that from 2004 dividend swaps were a major source of profit for multistrategy and macro hedge funds, because market-implied dividends traded at
high discounts while dividends grew. In this chapter's market, issuers' selling puts dividend futures 4\% below expected dividends for each year of
maturity. Buying the one-year future and holding it to expiry earns 4.2\% a year on average, but loses in one year in five, and 24\% in a recession
year. The same desk runs \omterm{def:s2:delta-one-and-dividends:indexarb}{index arbitrage} and \omterm{def:s2:delta-one-and-dividends:financing}{financing trades} on the index future. The build is \texttt{firm.deltaone}.

\section{Index arbitrage at a bank}

\begin{definition}[Index arbitrage]\label{def:s2:delta-one-and-dividends:indexarb}
\emph{Index arbitrage}\index{index arbitrage} is trading an index future against the basket of its constituents when the future's price departs from
its fair value (Book~1, chapter~21) by more than the cost of trading both, and unwinding when the gap closes.
\end{definition}

\texttt{firm.deltaone} simulates a year of minute-by-minute mispricings of the future around fair value, with a standard deviation of 4 basis points
of the index and a half-life of 10 minutes (\cref{lst:s2:delta-one-and-dividends:arb}). Trading the basket costs 2 basis points a side, for a bank that
crosses part of it internally (chapter~25), and the future 0.5, so a round trip costs 5 basis points. The desk sells the future and buys the basket when
the future is more than 5 basis points rich, the reverse when cheap, and unwinds when the gap closes:

\begin{center}
\small
\begin{tabular}{@{}lcccc@{}}
\toprule
execution lag & trades a year & bp per trade & bp a year & trades winning\\
\midrule
none & 2\,230 & 1.64 & 3\,657 & 100\%\\
1 minute & 2\,230 & 1.14 & 2\,539 & 69.2\%\\
2 minutes & 2\,230 & 0.77 & 1\,707 & 60.7\%\\
5 minutes & 2\,230 & $-0.26$ & $-584$ & 46.9\%\\
\bottomrule
\end{tabular}
\end{center}

Without a lag the trade cannot lose, since it enters beyond the cost and leaves at zero. Every minute of delay gives back part of a gap that closes in
ten, and at five minutes the trade loses money. \omterm{def:s2:delta-one-and-dividends:indexarb}{Index arbitrage} is a race whose prize is a basis point, and it belongs to whoever executes a basket
fastest and cheapest.

\section{Financing trades}

\begin{definition}[Financing trade]\label{def:s2:delta-one-and-dividends:financing}
A \emph{financing trade}\index{financing trade} is holding an asset against a derivative that delivers it later (a basket against a short future, or a
total-return swap in which the bank pays a client the asset's return), earning the financing spread implied by the derivative's price less the
bank's own cost of funding and of balance sheet.
\end{definition}

A future's price above spot, less dividends, implies a rate at which the market finances the basket. The synthetic implied rate is the overnight rate
plus a spread that averages about 20 basis points and reverts with a half-life of 40 days, plus 25 basis points in the last ten days of each quarter,
when banks' balance sheets are scarce (\cref{fig:s2:delta-one-and-dividends:spread}). The bank funds itself at 10 basis points over the overnight rate
and charges 5 for the balance sheet, a hurdle of 15. Buying the basket and selling a three-month future locks the implied spread for the quarter:

\begin{center}
\small
\begin{tabular}{@{}lccc@{}}
\toprule
entered each quarter & quarters taken & mean spread locked (bp) & bp of notional a year\\
\midrule
at the roll & 10\% & 15.4 & 0.4\\
in the quarter-end window & 67.5\% & 40.7 & 8.0\\
\bottomrule
\end{tabular}
\end{center}

At the roll, the spread rarely clears the hurdle and the round-trip cost. In the quarter-end window it usually does: the bank with balance sheet to spare
is paid by those who have none. A total-return swap is the same trade in another form, the bank holding the basket and paying the client its
return for a spread over funding.

\begin{omfigure}
\begin{tikzpicture}
\begin{axis}[width=11cm,height=5.2cm,xlabel={\small year},ylabel={\small implied spread over overnight (bp)},tick label style={font=\footnotesize},
  xmin=0,xmax=3,ymin=-25,ymax=65,grid=major,grid style={gray!25},table/col sep=comma,
  legend style={font=\scriptsize,at={(0.5,-0.3)},anchor=north},legend columns=2]
\addplot[omThm,thin] table[x=year,y=spread]{figdata/strategies-2/26-delta-one-and-dividends/spread.csv};
\addplot[omDef,thick,dashed] coordinates {(0,15) (3,15)};
\legend{futures-implied financing spread,bank's hurdle (funding and capital)}
\end{axis}
\end{tikzpicture}

\omcaption{Three years of the synthetic futures-implied financing spread over the overnight rate, with its quarter-end jumps, against the bank's own
funding and capital hurdle of 15 basis points. Data: \texttt{s2\_deltaone.market}.}\label{fig:s2:delta-one-and-dividends:spread}
\end{omfigure}

\section{Dividend risk from structured issuance}

An autocallable's buyer is in effect short a put on the index; the issuer holds the other side and hedges its delta by owning the index or its
future. Higher dividends lower the index's forward, raise the put's value and lower the note's value, a gain to the issuer, which is therefore long
dividends. Book~5's pricer (chapter~18) values a five-year autocallable with annual observations, an 8\% Phoenix coupon at a 70\% barrier and 60\%
protection at 101.84 per 100 at a 3\% dividend yield. A rise of 0.1 percentage point in the dividend yield lowers it by 0.11
(\cref{lst:s2:delta-one-and-dividends:div}): on a billion of notes, about 1.1 million for 10 basis points of dividend yield, over an expected life of
2.66 years.

\begin{definition}[Dividend supply pressure]\label{def:s2:delta-one-and-dividends:supply}
\emph{Dividend supply pressure}\index{dividend supply pressure} is the discount of dividend futures or swaps to expected dividends that arises when
hedgers who are naturally long dividends, such as issuers of structured products, sell more than other investors want to buy at fair value.
\end{definition}

\begin{dated}{2026-09}{Dividend futures}\label{dat:s2:delta-one-and-dividends:fexd}
Eurex's EURO STOXX 50 Index Dividend Futures (FEXD) pay EUR 100 per point of the index's dividends over the contract year, with expiries up to ten
years, cash settled on the third Friday of the expiry month. Eurex describes dividend derivatives as letting investors take positions in, or hedge,
future dividend payments, particularly for structured products and equity options.
\end{dated}

\section{Dividend trades}

The synthetic dividend market has expected dividends of 100.8 index points next year, growing 3\% a year, with 6\% of yearly noise and a 12\% chance
each year of a recession that cuts that year's dividends by 30\% and the index by 25\%. Futures trade at expected dividends less 4\% for each year of
maturity, the planted supply discount, so their term structure slopes down while expected dividends rise
(\cref{fig:s2:delta-one-and-dividends:dividends}). A buyer holding each future to expiry, over 20\,000 scenarios:

\begin{center}
\small
\begin{tabular}{@{}lccccc@{}}
\toprule
maturity (years) & 1 & 2 & 3 & 4 & 5\\
\midrule
expected dividends & 100.8 & 104.7 & 108.7 & 112.9 & 117.2\\
future & 96.7 & 96.3 & 95.7 & 94.9 & 93.7\\
mean return a year (\%) & 4.2 & 3.9 & 3.8 & 3.9 & 3.9\\
standard deviation a year (\%) & 12.3 & 9.1 & 7.5 & 6.6 & 5.9\\
5\% quantile a year (\%) & $-25.5$ & $-13.2$ & $-9.4$ & $-8.6$ & $-7.1$\\
chance of a loss (\%) & 20.9 & 25.2 & 29.8 & 28.0 & 24.6\\
\bottomrule
\end{tabular}
\end{center}

The one-year trade's return has a correlation of 0.41 with the index's return in the same year and loses 24.4\% on average in a recession year. The
dividend buyer is paid for supplying the insurance that the structured-product issuer does not want, and pays out in the same years as equities.
Van Binsbergen, Brandt and Koijen, recovering the prices of dividend strips, found that short-term dividend claims had higher expected returns, Sharpe
ratios and volatilities than the index, with CAPM betas below one. Kragt, de Jong and Driessen fitted the EURO STOXX 50 dividend futures curve with two
factors, one reverting within a year and one over the business cycle, and found that investors revise the value of dividends beyond the business
cycle only a little.

\begin{omfigure}
\begin{tikzpicture}
\begin{axis}[width=11cm,height=5cm,xlabel={\small maturity (years)},ylabel={\small index dividend points},tick label style={font=\footnotesize},
  xmin=0.8,xmax=5.2,ymin=90,ymax=120,xtick={1,2,3,4,5},grid=major,grid style={gray!25},table/col sep=comma,
  legend style={font=\scriptsize,at={(0.03,0.97)},anchor=north west}]
\addplot[black,thick,mark=*,mark size=1.2pt] table[x=year,y=expected]{figdata/strategies-2/26-delta-one-and-dividends/dividends.csv};
\addplot[omThm,thick,mark=square*,mark size=1.2pt] table[x=year,y=futures]{figdata/strategies-2/26-delta-one-and-dividends/dividends.csv};
\legend{expected dividends,dividend futures}
\end{axis}
\end{tikzpicture}

\omcaption{Expected dividends by year and the synthetic dividend futures curve, 4\% below them per year of maturity. Data:
\texttt{s2\_deltaone.dividends}.}\label{fig:s2:delta-one-and-dividends:dividends}
\end{omfigure}

\section{Strategy files}

\begin{strategyfile}{Index arbitrage}\label{strat:s2:delta-one-and-dividends:arb}
\sfield{Who pays you, and why} Traders who move the future ahead of the basket or the basket ahead of the future.
\sfield{Instruments and venues} Index futures; the constituent basket; ETFs.
\sfield{Signal} The future against fair value with the market's financing and dividends.
\sfield{Sizing and execution} Beyond the round-trip cost; the basket executed in one program.
\sfield{Costs} Basket and futures spreads; latency.
\sfield{How it dies} Faster competitors; a wrong dividend or financing input.
\sfield{Horizon, capacity, infrastructure} Minutes; basket execution and low latency.
\sfield{Backtest honestly} Execute after the signal, at the prices then available.
\sfield{Sources} This chapter: 1.14 basis points per trade at a one-minute lag, a loss at five minutes.
\end{strategyfile}

\begin{strategyfile}{Total-return swap financing}\label{strat:s2:delta-one-and-dividends:trs}
\sfield{Who pays you, and why} Investors who want the index's return without funding it; scarce balance sheet at quarter-ends.
\sfield{Instruments and venues} Total-return swaps; basket against futures.
\sfield{Signal} The implied financing spread against the bank's funding and capital charge.
\sfield{Sizing and execution} Within balance-sheet limits; more in the quarter-end window.
\sfield{Costs} Funding; capital; basket trading.
\sfield{How it dies} Regulation that raises the capital charge; competition for balance sheet.
\sfield{Horizon, capacity, infrastructure} Months; treasury and balance-sheet allocation.
\sfield{Backtest honestly} The bank's own funding and capital costs at the time.
\sfield{Sources} This chapter: 8.0 basis points of notional a year entered at quarter-end.
\end{strategyfile}

\begin{strategyfile}{Long dividends against structured supply}\label{strat:s2:delta-one-and-dividends:dividends}
\sfield{Who pays you, and why} Issuers of structured products selling the dividends their hedges leave them with.
\sfield{Instruments and venues} Index dividend futures and swaps.
\sfield{Signal} The discount of the futures to forecast dividends.
\sfield{Sizing and execution} Held to expiry; sized for a recession cut.
\sfield{Costs} Futures spreads; margin.
\sfield{How it dies} Dividend cuts in recessions, including regulatory bans on bank dividends.
\sfield{Horizon, capacity, infrastructure} One to five years.
\sfield{Backtest honestly} Include recession years; forecast dividends as of the trade date.
\sfield{Sources} Manley and Mueller-Glissmann (2008); van Binsbergen, Brandt and Koijen (2012); this chapter: 4.2\% a year, $-24.4\%$ in a recession
year.
\end{strategyfile}

\begin{strategyfile}{Dividend calendar spread}\label{strat:s2:delta-one-and-dividends:calendar}
\sfield{Who pays you, and why} Supply concentrated at the maturities structured products need.
\sfield{Instruments and venues} Dividend futures of two maturities.
\sfield{Signal} The curve's slope against expected dividend growth.
\sfield{Sizing and execution} Weighted so that a parallel change in dividends nets out.
\sfield{Costs} Two legs.
\sfield{How it dies} A recession that cuts near dividends more than far ones.
\sfield{Horizon, capacity, infrastructure} Years.
\sfield{Backtest honestly} A model of the term structure fitted without the future.
\sfield{Sources} Kragt, de Jong and Driessen (2020); no performance figure verified.
\end{strategyfile}

\section{Tutorial: dividends nobody wanted}\label{tut:s2:delta-one-and-dividends:tut}

\begin{tutorial}
\textbf{Goal.} Run \omterm{def:s2:delta-one-and-dividends:indexarb}{index arbitrage} with execution lags, \omterm{def:s2:delta-one-and-dividends:financing}{financing trades} at the roll and at quarter-end, the issuer's dividend exposure from an
autocallable, and a dividend market with supply pressure. \textbf{End state:} the four tables and two figures.
\begin{tutsteps}
\item \textbf{Arbitrage and financing}.
\omcode{code/firm/deltaone/firm_deltaone.py}{78}{104}{Fading mispricings after a lag; locking the implied financing spread.}\label{lst:s2:delta-one-and-dividends:arb}
\item \textbf{Dividends}.
\omcode{code/firm/deltaone/firm_deltaone.py}{107}{136}{The issuer's dividend exposure; dividend futures with a supply discount.}\label{lst:s2:delta-one-and-dividends:div}
\item \textbf{Run} \texttt{arbitrage()}, \texttt{financing()}, \texttt{exposure()}, \texttt{dividends()} and \texttt{fig\_deltaone.py}.
\end{tutsteps}
\textbf{What to change next.} Let the supply discount rise with issuance and fall after a recession; hedge the dividend trade with index puts; add a
dividend-ban scenario for banks.
\end{tutorial}

\section{Build: delta-one}\label{bld:s2:delta-one-and-dividends:deltaone}

\begin{build}
\bfield{Purpose} \omterm{def:s2:delta-one-and-dividends:indexarb}{Index arbitrage}, \omterm{def:s2:delta-one-and-dividends:financing}{financing trades}, the issuer's dividend exposure and a dividend market with supply pressure.
\bfield{Interface} \texttt{DeltaOneConfig(\dots)}, \texttt{simulate\_financing(cfg)}, \texttt{index\_arb(sim, cfg, lag)},
\texttt{financing\_trade(sim, cfg, offset)}, \texttt{issuer\_dividend\_exposure(vol, r, q, bump)}, \texttt{dividend\_market(cfg)}.
\bfield{Rules} Arbitrage beyond the round-trip cost, executed after a lag; \omterm{def:s2:delta-one-and-dividends:financing}{financing trades} only when they clear the hurdle; futures at expected
dividends less the discount.
\bfield{Acceptance tests} \texttt{code/firm/deltaone/tests/}: no lag, no loss; the \omterm{def:s2:delta-one-and-dividends:financing}{financing trade} by hand; the issuer is long dividends; no discount,
no expected return.
\bfield{Stretch} Dividend swaps with single stocks; supply that responds to issuance; stochastic interest rates.
\end{build}

\begin{omsources}
\item J.~van Binsbergen, M.~Brandt and R.~Koijen, ``On the timing and pricing of dividends'', \emph{American Economic Review} 102(4), 2012.
\item R.~Manley and C.~Mueller-Glissmann, ``The market for dividends and related investment strategies'', \emph{Financial Analysts Journal} 64(3),
2008.
\item J.~Kragt, F.~de Jong and J.~Driessen, ``The dividend term structure'', \emph{Journal of Financial and Quantitative Analysis} 55(3), 2020.
\item Eurex, ``Dividend derivatives'', factsheet, accessed 25 September 2026.
\end{omsources}

\section{Exercises}

\begin{exercise}[$\star$]\label{exo:s2:delta-one-and-dividends:1}
The future is 7 basis points rich and a round trip costs 5. The gap closes completely. What does the arbitrage earn on \$50 million?
\end{exercise}

\begin{exercise}[$\star$]\label{exo:s2:delta-one-and-dividends:2}
The implied spread is 40 basis points, the bank's hurdle 15, and a quarter's round trip costs 5 basis points. What does a quarter's \omterm{def:s2:delta-one-and-dividends:financing}{financing trade}
earn on \$100 million?
\end{exercise}

\begin{exercise}[$\star$]\label{exo:s2:delta-one-and-dividends:3}
Expected dividends are 108.7 points and the future trades at 95.7. What does holding it to expiry earn if dividends come in as expected, in total and a
year over three years?
\end{exercise}

\begin{exercise}[$\star\star$]\label{exo:s2:delta-one-and-dividends:4}
Why is the issuer of an autocallable long dividends?
\end{exercise}

\begin{exercise}[$\star\star$]\label{exo:s2:delta-one-and-dividends:5}
Short-term dividend claims earned more than the index in van Binsbergen, Brandt and Koijen. Why might supply explain part of it?
\end{exercise}

\begin{exercise}[$\star\star$]\label{exo:s2:delta-one-and-dividends:6}
Why does the implied financing spread jump at quarter-ends?
\end{exercise}

\begin{exercise}[$\star\star\star$]\label{exo:s2:delta-one-and-dividends:7}
\emph{Coding.} Run \texttt{index\_arb} with lags of 1, 2 and 5 minutes. At what lag does the trade stop paying, and why?
\end{exercise}

\begin{exercise}[$\star\star\star$]\label{exo:s2:delta-one-and-dividends:8}
\emph{Find the flaw.} ``Dividend futures have earned about 4\% a year with low volatility; they are a bond substitute.''
\end{exercise}

\section{Problem: Dividends Nobody Wanted}

\begin{problem}[{Weekend problem --- delta-one and dividends}]\label{pb:s2:delta-one-and-dividends:1}
The chapter's synthetic desk and the public record.

\medskip\noindent\textbf{Part I --- Arbitrage.}
\begin{enumerate}
\item Define \omterm{def:s2:delta-one-and-dividends:indexarb}{index arbitrage}.
\item Describe the mispricing and the costs.
\item Give the results by execution lag.
\item Why can the trade not lose without a lag?
\end{enumerate}

\medskip\noindent\textbf{Part II --- Financing.}
\begin{enumerate}[resume]
\item Define a \omterm{def:s2:delta-one-and-dividends:financing}{financing trade}.
\item What does the future's price imply about financing?
\item Give the results at the roll and at quarter-end.
\item How is a total-return swap the same trade?
\end{enumerate}

\medskip\noindent\textbf{Part III --- Dividends.}
\begin{enumerate}[resume]
\item Why is the autocallable's issuer long dividends, and by how much?
\item Define \omterm{def:s2:delta-one-and-dividends:supply}{dividend supply pressure}.
\item What does the Eurex contract pay?
\item Give the dividend table.
\end{enumerate}

\medskip\noindent\textbf{Part IV --- The verdict.}
\begin{enumerate}[resume]
\item State the \emph{named result}: the dividend trade's return from the supply discount.
\item What did Manley and Mueller-Glissmann report?
\item What did van Binsbergen, Brandt and Koijen find?
\item What did Kragt, de Jong and Driessen find?
\item When does the dividend trade lose, and why then?
\item Which strategy file needs balance sheet?
\item How does this chapter relate to chapter~12's swap spreads?
\item In one sentence: what does a delta-one desk sell?
\end{enumerate}
\end{problem}

\section{Interview questions}

\begin{interviewq}[$\star$ \iqroles{trader}]\label{iq:s2:delta-one-and-dividends:1}
How do you compute an index future's fair value, and what inputs can be wrong?
\end{interviewq}

\begin{interviewq}[$\star\star$ \iqroles{researcher}]\label{iq:s2:delta-one-and-dividends:2}
How would you forecast next year's index dividends?
\end{interviewq}

\begin{interviewq}[$\star\star$ \iqroles{trader}]\label{iq:s2:delta-one-and-dividends:3}
A client wants a three-month total-return swap on an index over quarter-end. How do you price it?
\end{interviewq}

\begin{interviewq}[$\star\star$ \iqroles{risk}]\label{iq:s2:delta-one-and-dividends:4}
What stress would you apply to a book long dividend futures?
\end{interviewq}

\begin{interviewq}[$\star\star$ \iqroles{developer}]\label{iq:s2:delta-one-and-dividends:5}
Design the system that keeps index fair values up to date intraday.
\end{interviewq}

\begin{interviewq}[$\star\star\star$ \iqroles{researcher}]\label{iq:s2:delta-one-and-dividends:6}
The future's price is $F = (S - D) e^{rT}$ with $D$ the present value of dividends to expiry. Show that an error $\delta$ in $D$ shifts fair value by
$-\delta e^{rT}$, and say why a dividend forecast error looks like an arbitrage.
\end{interviewq}
